% ============================================================
%  Manuscript for European Journal of Nuclear Medicine and
%  Molecular Imaging (EJNMMI) — Original Article
% ============================================================

\documentclass[sn-basic]{sn-jnl}

\usepackage{graphicx}
\usepackage{booktabs}
\usepackage{multirow}
\usepackage{threeparttable}
\usepackage{hyperref}
\usepackage{amsmath}
\usepackage{siunitx}
\usepackage[utf8]{inputenc}
\usepackage[numbers,sort&compress]{natbib}

% ============================================================
% Article meta-data
% ============================================================
\title[VLM-based diagnosis with automatic ROI segmentation]{%
  Vision-Language Model–Based Diagnosis of Renal Dynamic
  Scintigraphy with UNet+LTAE Automatic Kidney Segmentation:
  An Ablation Study}

\author*[1]{\fnm{Xiaoxiao} \sur{Zhang}}
\author[1]{\fnm{Anonymous} \sur{Author}}
\author[1]{\fnm{Anonymous} \sur{Author}}

\affil*[1]{%
  \org{Department of Nuclear Medicine, XXX Hospital}, %
  \org{\city{XXX City}, \country{China}}%
}

\correspondence{%
  \email{xxx@hospital.edu.cn}
}

% ============================================================
\begin{document}

\maketitle

% ============================================================
\section*{Abstract}
% ============================================================

\textbf{Purpose:}
Manual region-of-interest (ROI) delineation in renal dynamic
scintigraphy is operator-dependent and time-consuming.  We
evaluate whether an automatic deep-learning segmentation
pipeline (UNet with Lightweight Triplet Attention and
Efficient channel adaptation, UNet+LTAE) can replace
physician-drawn ROIs in a vision-language model (VLM)–based
diagnostic workflow for 99mTc-EC diuretic renography, and
quantify the contribution of time-activity curves (TACs) and
image features via three orthogonal ablations.

\textbf{Methods:}
130-frame 99mTc-EC dynamic sequences from 60 patients (20
functional, 20 mechanical, 20 mixed obstruction) were fused
into 26 temporal frames.  A UNet+LTAE model was trained on
102 samples from an independent cohort (75/25 split, 200
epochs, BCE+Dice loss) to segment kidney ROIs automatically
from fused images.  A two-stage VLM pipeline using GPT-5.5
(temperature 0.0) first extracted 11 structured evidence
fields from kidney crops and TAC curves (Stage~1), then
synthesized a four-class diagnosis (normal, functional
obstruction, mechanical obstruction, mixed obstruction)
(Stage~2).  Three ablations were performed on 52 cases (104
kidneys):~(i)~removing TAC input,~(ii)~removing image
input,~(iii)~replacing physician ROIs with UNet+LTAE
segmentations.  Accuracy, Cohen's~$\kappa$,
quadratic-weighted~$\kappa$, per-class~F1 with 95\%
bootstrap confidence intervals, and McNemar's test were
computed.

\textbf{Results:}
The full pipeline with physician ROIs achieved 80.77\%
accuracy ($\kappa=0.730$).  UNet+LTAE automatic ROIs yielded
78.85\% ($\kappa=0.703$), without significant difference
(McNemar $p=0.625$).  Removing TACs severely degraded
performance to 42.31\% ($\kappa=0.240$, $p<$0.001), whereas
TAC-only input achieved 67.31\% ($\kappa=0.552$,
$p=0.0013$).  Mechanical obstruction was best identified
(F1~0.889), normal kidneys showed robust discrimination
(F1~0.920), while mixed obstruction remained the most
challenging category (F1~0.545).

\textbf{Conclusion:}
UNet+LTAE automatic kidney segmentation can replace manual
ROI delineation without significant diagnostic accuracy
loss.  TAC curves are the dominant information source, yet
image features provide complementary value.  The two-stage
VLM pipeline offers explainable, structured reasoning and
achieves clinically competitive diagnostic accuracy.

\vspace{1em}
\noindent\textbf{Keywords:}
renal dynamic scintigraphy, vision-language model, deep
learning, UNet, automatic segmentation, time-activity curve,
diuretic renography, obstruction classification

% ============================================================
\section{Introduction}
\label{sec:intro}
% ============================================================

Diuretic renography with \textsuperscript{99m}Tc-EC is a
first-line imaging modality for evaluating suspected renal
outflow obstruction~\cite{oneill,rossleigh}.  The standard
interpretation integrates qualitative assessment of cortical
uptake, parenchymal transit, and drainage patterns with
quantitative time-activity curves (TACs) derived from
manually drawn regions of interest (ROIs) over each
kidney~\cite{gordon,taylor}.  Despite its established
clinical utility, the procedure relies heavily on operator
expertise: manual ROI delineation is time-consuming, shows
moderate inter-observer variability~\cite{degaard}, and must
be performed before any quantitative analysis can proceed.

Deep learning has substantially advanced medical image
segmentation, with U-Net architectures dominating
applications in renal imaging~\cite{ronneberger,
isensee}.  However, most prior work focuses on static
anatomical segmentation (e.g.,~CT or MR
kidney/parenchyma boundaries) rather than dynamic
scintigraphic sequences where the target structures are
poorly defined, low-count, and vary in contrast throughout
the acquisition.  To our knowledge, automated ROI
delineation for TAC extraction in renal dynamic scintigraphy
followed by downstream diagnostic reasoning has not been
systematically evaluated.

Parallel advances in vision-language models
(VLMs)~\cite{openai,gpt4} offer a new paradigm for
diagnostic image interpretation.  Unlike traditional
deep-learning classifiers that output a label without
explanation, VLMs can generate structured, interpretable
evidence chains that mimic the radiologist's reasoning
process.  Recent studies have applied VLMs to
chest~X-ray~\cite{wu}, retinal~photography~\cite{luo}, and
general nuclear medicine
reporting~\cite{sibille,werner}, but a dedicated evaluation
for renal dynamic scintigraphy with comprehensive ablation
analysis is lacking.

In this study, we present and evaluate a complete pipeline
for renal dynamic scintigraphy diagnosis that combines:
(1)~a UNet with Lightweight Triplet Attention and Efficient
channel adaptation (UNet+LTAE) for automatic kidney ROI
segmentation, (2)~a two-stage VLM-based diagnostic workflow
(Stage~1: structured image and TAC evidence extraction;
Stage~2: evidence synthesis into final diagnosis), and
(3)~three orthogonal ablation experiments that quantify the
contribution of each pipeline component: TAC curves, image
features, and automatic versus manual ROIs.  Our goal is to
determine whether automatic segmentation can replace manual
delineation without compromising diagnostic accuracy, and to
characterize the relative importance of image-derived versus
TAC-derived information in VLM-based renography
interpretation.

% ============================================================
\section{Materials and Methods}
\label{sec:methods}
% ============================================================

\subsection{Dataset and Preprocessing}

We retrospectively collected 99mTc-EC diuretic renography
studies from 60 patients (mean age~54~years, 34~male) who
underwent standard diuretic renography at our institution
between January 2023 and December 2023.  The cohort
comprised three groups based on final clinical diagnosis:
function obstruction (n=20), mechanical obstruction (n=20),
and mixed obstruction (n=20).  Normal kidneys (contralateral
unaffected kidneys within the same cohort) were labeled as
having no obstruction.  Ground truth diagnoses were
established by consensus of two experienced nuclear medicine
physicians (15+~years) based on clinical history, imaging
findings, and, where available, surgical outcome or
long-term follow-up.

Each study consisted of 130~dynamic frames (128~$\times$~128
matrix, 2~s per frame for the first 2~min, followed by
longer intervals for the remaining~8~min).  Diuretic
(40~mg~furosemide) was administered intravenously at
approximately frame~85, consistent with the F+10
protocol~\cite{official}.  The raw 130-frame sequences were
temporally fused into 26~time steps following a standard
protocol: the first 30~frames were averaged into a single
frame, frames~31--126 were grouped in blocks of~4, and the
final 2~frames were averaged.  This fusion reduces image
noise while preserving the essential temporal dynamics.

For each case, kidney crops were generated from a subset of
11~representative frames (indices~0,3,6,9,12,14,16,18,21,24,25
of the 26 fused images) at 512~$\times$~512 resolution with
contrast enhancement (factor~2.2) and 0.4~padding, using ROI
boundaries from either physician-drawn polygons or
model-predicted masks.

Of the 60~patients, 52~(104~kidneys) had complete pipeline
results available for all four experimental conditions and
constitute the primary evaluation set.  The remaining
8~patients in the mixed-obstruction group were excluded due
to missing segmentation checkpoints for that specific cohort
subgroup.

\subsection{UNet+LTAE for Automatic Kidney Segmentation}

\begin{figure}[t]
  \centering
  \includegraphics[width=\textwidth]{fig_unet_architecture.pdf}
  \caption{Architecture of the UNet+LTAE model.  The network
    takes 26-channel fused images as input and produces a
    binary kidney segmentation mask.  LTAE~blocks apply
    channel-wise attention via adaptive average pooling,
    a~two-layer MLP (reduction ratio~4), and sigmoid gating
    with residual connection.}
  \label{fig:unet}
\end{figure}

The segmentation model (Fig.~\ref{fig:unet}) adopts a 4-level
encoder---decoder U-Net backbone~\cite{ronneberger}.  Each
encoder level consists of two consecutive
Conv2D--BatchNorm2D--ReLU operations (DoubleConv),
followed by 2$\times$2 max pooling with stride~2.  Channel
dimensions progress from~64 to~128,~256,~512, with a
1024-channel bottleneck.  The decoder mirrors the encoder
using transposed convolutions and skip connections.  The
LTAE~Block---Lightweight Triplet Attention and Efficient
channel adaptation---is inserted at the bottleneck and
computes channel attention via
AdaptiveAvgPool2D(1)~$\rightarrow$~MLP(reduction=4)
~$\rightarrow$~Sigmoid, applied as a residual multiplication:
$\mathbf{x}' = \mathbf{x} \odot \sigma(\text{MLP}(\text{GAP}(\mathbf{x}))) + \mathbf{x}$.
The final layer is a 1$\times$1 convolution producing a
single-channel output, followed by sigmoid activation.

The model was trained on 102~samples from an independent
cohort (data3) with a 75/25~random split (76 train, 26~val,
seed~42).  Inputs were the 26-channel fused images at
256~$\times$~256 resolution; targets were physician-drawn
polygon masks rasterized at the same resolution.  Loss
function: $\mathcal{L} = 0.5 \cdot
\text{BCEWithLogitsLoss} + 0.5 \cdot \text{DiceLoss}$.
Optimizer: Adam (lr~=~10$^{-4}$), batch size~2, 200~epochs
with ReduceLROnPlateau (factor~0.5, patience~5).  Data
augmentation included random rotation ($\pm$15$^\circ$),
horizontal flip, vertical flip, and ElasticTransform.
Best model was selected by validation Dice score.

At inference, the predicted binary mask is post-processed
using connected-component analysis.  When a single component
is detected, left/right kidney separation is performed by
splitting at the center-of-mass~x coordinate; when multiple
components exist, the two largest are assigned to left and
right kidneys by their~x centroids.

\subsection{Time-Activity Curve Extraction}

TACs were derived by computing the mean pixel intensity
within each kidney ROI across all 130~raw frames.  Since the
raw images are inverted (higher counts appear darker), we
applied 255-inversion before averaging.  The resulting
curves were normalized to unit maximum and smoothed with a
5-frame moving-average filter.  TACs were visualized as
plot images (intensity vs.\ frame index) and provided as
input to the VLM alongside the kidney crop images.

\subsection{Two-Stage VLM Diagnostic Pipeline}

\begin{figure}[t]
  \centering
  \includegraphics[width=\textwidth]{fig_pipeline.pdf}
  \caption{Two-stage diagnostic pipeline.  Stage~1: a VLM
    (GPT-5.5, temperature~0.0) receives 11~kidney crop images
    and TAC plot images and extracts 11~structured evidence
    fields with controlled vocabulary.  Stage~2: an LLM
    (same model) synthesizes the evidence into a four-class
    diagnosis with confidence and reasoning.}
  \label{fig:pipeline}
\end{figure}

The diagnostic pipeline (Fig.~\ref{fig:pipeline}) operates in
two stages using GPT-5.5 with temperature~0.0 across all
calls for deterministic output.

\paragraph{Stage~1: Structured Evidence Extraction}
For each kidney, the VLM simultaneously receives 11~cropped
frames (fused images at selected time points) and the
corresponding TAC plot image.  The prompt (Appendix~A)
instructs the model to extract 11~evidence fields using a
strict controlled vocabulary:
\begin{enumerate}
  \item \textit{Early visualization} (timely / mildly delayed / severely delayed)
  \item \textit{Peak timing} (early / mid / late peaking)
  \item \textit{Pre-diuretic trend} (declining / plateau / rising)
  \item \textit{Post-diuretic drop} (present / absent)
  \item \textit{Post-diuretic emptying} (sufficient / partial / insufficient)
  \item \textit{Collecting system change} (clearly reduced / partially reduced / unchanged)
  \item \textit{Post-diuretic response} (clear improvement / some improvement / no improvement)
  \item \textit{Late-phase retention} (low / moderate / high)
  \item \textit{Late retention relative to diuretic} (lower than / similar to / higher than)
  \item \textit{Overall dynamic pattern} (coordinated emptying / delayed emptying / persistent retention)
  \item \textit{Reasoning focus} (free-text summary)
\end{enumerate}
The output is validated against the allowed vocabulary; any
field with an unrecognized value is replaced by
``uncertain.''

\paragraph{Stage~2: Evidence Synthesis}
The Stage~1 structured outputs for both kidneys are passed to
a second LLM call that synthesizes the evidence into a final
diagnosis for each kidney (normal / functional obstruction /
mechanical obstruction / mixed obstruction), along with a
confidence level (low / moderate / high) and a brief
reasoning statement.  The model also provides an overall
assessment of image-TAC consistency and overall confidence.

\subsection{Ablation Study Design}

We designed three orthogonal ablation experiments to isolate
the contribution of each pipeline component:

\begin{enumerate}
  \item[\textbf{A1}] \textbf{No-TAC}: Stage~1 receives only
    the 11~kidney crop images without TAC curves.  This
    isolates the diagnostic value of spatial/morphological
    image features alone.
  \item[\textbf{A2}] \textbf{TAC-only}: Stage~1 receives only
    the TAC plot images without kidney crops.  This isolates
    the diagnostic value of temporal uptake/washout dynamics.
  \item[\textbf{A3}] \textbf{UNet+LTAE vs.\ physician ROI}:
    The complete pipeline is run identically except that
    kidney crops and TACs are derived from UNet+LTAE
    automatic segmentations instead of physician-drawn
    polygons.  This evaluates whether automatic segmentation
    can substitute for manual delineation.
\end{enumerate}

All four conditions (physician ROI full pipeline + three
ablations) were evaluated on the same set of 52~cases (104
kidneys) to enable paired statistical comparison.

\subsection{Statistical Analysis}

Performance was assessed using overall accuracy,
Cohen's~$\kappa$ (linear-weighted), and quadratic-weighted
$\kappa$.  Per-class performance was characterized by
sensitivity, specificity, positive predictive value (PPV),
negative predictive value (NPV), and F1~score.  All metrics
were computed with 95\%~bootstrap confidence intervals (1000
iterations, random seed~42).  Pairwise comparisons between
pipeline configurations were performed using McNemar's test
for paired nominal data.  A~$p$-value~$<$~0.05 was
considered statistically significant.  The label mapping
``normal kidney'' (model output) $\rightarrow$ ``normal
obstruction'' (ground truth) was applied as a synonym
normalization step.

% ============================================================
\section{Results}
\label{sec:results}
% ============================================================

\subsection{Overall Diagnostic Performance}

Table~\ref{tab:overall} summarizes the diagnostic performance
of all four pipeline configurations on the 52-case (104
kidney) evaluation set.  The full pipeline with physician
ROIs achieved 80.77\%~accuracy ($\kappa=0.730$), serving as
the reference standard for comparison.  UNet+LTAE automatic
ROIs yielded 78.85\%~accuracy ($\kappa=0.703$), a
non-significant difference of 1.92~percentage points
(McNemar~$p=0.625$).  Removing TAC input caused a dramatic
decline to 42.31\%~accuracy ($\kappa=0.240$), significantly
worse than the full pipeline ($p<0.001$).  TAC-only input
achieved 67.31\%~accuracy ($\kappa=0.552$), significantly
lower than the full pipeline ($p=0.0013$) but substantially
higher than the no-TAC condition.

\begin{table}[t]
  \centering
  \caption{Overall diagnostic performance across four
    pipeline configurations (52~cases, 104~kidneys).  95\%
    bootstrap confidence intervals are shown in parentheses.}
  \label{tab:overall}
  \small
  \begin{threeparttable}
    \begin{tabular}{lcccc}
      \toprule
      Method & Accuracy (\%) & Cohen's~$\kappa$ & QW~$\kappa$ & Macro~F1 \\
      \midrule
      Physician ROI (full) & 80.77 (72.15--87.19) & 0.730 (0.623--0.825) & 0.868 & 0.763 \\
      UNet+LTAE & 78.85 (70.04--85.59) & 0.703 (0.594--0.800) & 0.859 & 0.741 \\
      No TAC & 42.31 (33.25--51.91) & 0.240 (0.126--0.361) & 0.577 & 0.418 \\
      TAC only & 67.31 (57.82--75.57) & 0.552 (0.436--0.664) & 0.777 & 0.614 \\
      \bottomrule
    \end{tabular}
  \end{threeparttable}
\end{table}

\subsection{Per-Class Performance}

Per-class metrics for the physician-ROI and UNet+LTAE
pipelines are presented in Table~\ref{tab:perclass}.  Both
configurations showed consistent patterns: mechanical
obstruction was identified with high sensitivity and
specificity (F1~0.889 and~0.864, respectively), and normal
kidneys were reliably distinguished from pathological cases
(F1~0.920 for both).  Functional obstruction showed moderate
discriminative performance (F1~0.698 and~0.667).  Mixed
obstruction was the most challenging category, with
sensitivity of~56.3\% and F1~scores of~0.546 and~0.514 for
physician and UNet+LTAE pipelines, respectively.

\begin{table}[t]
  \centering
  \caption{Per-class performance for physician-ROI and
    UNet+LTAE pipelines (52~cases, 104~kidneys).  95\%
    bootstrap confidence intervals in parentheses.}
  \label{tab:perclass}
  \small
  \begin{threeparttable}
    \begin{tabular}{lcccc}
      \toprule
      Class & Sensitivity & Specificity & PPV & F1 \\
      \midrule
      \multicolumn{5}{l}{\textit{Physician ROI}} \\
      \hspace{1em} Normal kidney & 88.89 (76.5--95.2) & 96.61 (88.5--99.1) & 95.24 & 0.920 (0.854--0.970) \\
      \hspace{1em} Functional & 65.22 (44.9--81.2) & 93.83 (86.4--97.3) & 75.00 & 0.698 (0.513--0.837) \\
      \hspace{1em} Mixed & 56.25 (33.2--76.9) & 90.91 (83.1--95.3) & 52.94 & 0.546 (0.308--0.732) \\
      \hspace{1em} Mechanical & 100.0 (83.9--100) & 94.05 (86.8--97.4) & 80.00 & 0.889 (0.776--0.976) \\
      \midrule
      \multicolumn{5}{l}{\textit{UNet+LTAE}} \\
      \hspace{1em} Normal kidney & 88.89 (76.5--95.2) & 96.61 (88.5--99.1) & 95.24 & 0.920 (0.849--0.972) \\
      \hspace{1em} Functional & 60.87 (40.8--77.8) & 93.83 (86.4--97.3) & 73.68 & 0.667 (0.476--0.813) \\
      \hspace{1em} Mixed & 56.25 (33.2--76.9) & 88.64 (80.3--93.7) & 47.37 & 0.514 (0.286--0.698) \\
      \hspace{1em} Mechanical & 95.00 (76.4--99.1) & 94.05 (86.8--97.4) & 79.17 & 0.864 (0.737--0.957) \\
      \bottomrule
    \end{tabular}
  \end{threeparttable}
\end{table}

\subsection{Ablation Analysis}

Table~\ref{tab:mcnemar} presents pairwise McNemar comparisons
among the four pipeline configurations.  The key findings
are:

\begin{enumerate}
  \item \textbf{UNet+LTAE vs.\ physician ROI}: No significant
    difference ($p=0.625$), supporting the feasibility of
    replacing manual ROIs with automatic segmentations.
  \item \textbf{Full pipeline vs.\ no TAC}: Highly significant
    difference (both $p<0.001$), confirming that TAC curves
    are the dominant diagnostic information source.
  \item \textbf{Full pipeline vs.\ TAC-only}: Significant
    difference ($p=0.0013$ and $p=0.0075$ for physician and
    UNet+LTAE respectively), demonstrating that image
    features provide complementary diagnostic value beyond
    TAC curves alone.
  \item \textbf{TAC-only vs.\ no TAC}: Significant difference
    ($p=0.0012$), showing that temporal dynamics alone
    outperform spatial features alone.
\end{enumerate}

\begin{table}[t]
  \centering
  \caption{McNemar~$p$-values for pairwise comparisons among
    pipeline configurations.}
  \label{tab:mcnemar}
  \small
  \begin{threeparttable}
    \begin{tabular}{lcccc}
      \toprule
      & Physician ROI & UNet+LTAE & No TAC & TAC only \\
      \midrule
      Physician ROI & --- & 0.625 & $<$0.001 & 0.0013 \\
      UNet+LTAE & 0.625 & --- & $<$0.001 & 0.0075 \\
      No TAC & $<$0.001 & $<$0.001 & --- & 0.0012 \\
      TAC only & 0.0013 & 0.0075 & 0.0012 & --- \\
      \bottomrule
    \end{tabular}
  \end{threeparttable}
\end{table}

\subsection{Confusion Matrix Analysis}

Figure~\ref{fig:confusion} presents the confusion matrices
for the physician-ROI and UNet+LTAE pipelines.  Both
configurations show a similar error pattern.  Mixed
obstruction was most frequently misclassified as functional
obstruction or mechanical obstruction.  Functional
obstruction errors were predominantly toward normal and
mixed categories.  Mechanical obstruction showed the most
consistent classification, with near-perfect sensitivity in
the physician-ROI configuration and only a single
misclassification in the UNet+LTAE configuration.

\begin{figure}[t]
  \centering
  \includegraphics[width=0.85\textwidth]{fig_confusion.pdf}
  \caption{Confusion matrices for physician-ROI (left) and
    UNet+LTAE (right) pipelines.  Rows: ground truth;
    columns: predicted labels.  Values indicate number of
    kidneys.}
  \label{fig:confusion}
\end{figure}

% ============================================================
\section{Discussion}
\label{sec:discussion}
% ============================================================

This study demonstrates three principal findings.  First,
UNet+LTAE automatic kidney segmentation achieves diagnostic
accuracy comparable to physician-drawn ROIs (78.85\% vs.\
80.77\%, $p=0.625$), indicating that the manual ROI step can
be reliably automated in a VLM-based diagnostic pipeline.
Second, the ablation experiments reveal that TAC curves
provide the dominant contribution to diagnostic performance:
removing them causes accuracy to nearly halve (42.31\%),
while TAC-only input alone achieves 67.31\%.  Third, image
features provide complementary information beyond TAC curves,
as the full pipeline significantly outperforms TAC-only input
($p=0.0013$), consistent with the clinical practice of
integrating both visual pattern assessment and quantitative
curve analysis~\cite{taylor,gordon}.

The comparable performance between automatic and manual ROIs
is particularly noteworthy.  The 1.92~percentage-point
difference is not statistically significant, and the
per-class F1~scores remain stable across all four diagnostic
categories.  This suggests that, at least for the VLM-based
interpretation framework, small differences in ROI boundary
placement do not materially alter the extracted TAC morphology
or the spatial features captured by the kidney crop images.
This robustness has practical implications: it removes a
major barrier to fully automated renography reporting, as
manual ROI delineation is both the most time-consuming step
and a source of inter-operator variability~\cite{degaard}.

The ablation results provide quantitative evidence for
longstanding clinical intuition.  The 38.46~percentage-point
drop when TACs are removed confirms that temporal dynamics---
the uptake, peak, and washout pattern---are the cornerstone
of diuretic renography interpretation.  Conversely, the
persistent gap between TAC-only (67.31\%) and the full
pipeline (80.77\%) demonstrates that image features (renal
cortical uptake pattern, collecting system morphology, and
their evolution over time) contribute meaningful independent
information.  This aligns with established
protocols~\cite{official} that recommend assessing both
quantitative curve parameters and qualitative image features.

The per-class analysis reveals a consistent difficulty
gradient across both automatic and manual configurations.
Mechanical obstruction and normal kidneys are classified
with high confidence, reflecting their distinctive TAC
phenotypes: mechanical obstruction shows the classic rising
or plateau pattern with absent diuretic response, while
normal kidneys exhibit prompt washout.  Functional
obstruction, characterized by delayed but ultimately
complete drainage, has intermediate features that overlap
with both normal and mechanical patterns.  Mixed obstruction
presents the greatest diagnostic challenge, as it combines
mechanical and functional components whose relative
contributions vary across cases, creating a continuous
spectrum rather than discrete categories.

\subsection{Limitations}

Several limitations should be acknowledged.  First, the
dataset size (52~cases, 104~kidneys) is modest, and although
bootstrap confidence intervals provide some measure of
precision, the estimates for rarer classes (particularly
functional and mixed obstruction) have wide intervals.
Second, the UNet+LTAE model was trained on 26-channel fused
images; fine-grained temporal information available in the
full 130-frame sequence may improve segmentation accuracy.
Third, all VLM evaluations used a single model (GPT-5.5);
generalizability across different foundation models warrants
investigation.  Fourth, the 8~excluded mixed-obstruction
cases represent a potential selection bias, as the most
severe mixed cases may have been preferentially excluded.
Finally, the study is retrospective and does not evaluate
the pipeline's impact on clinical workflow efficiency or
reader confidence.

\subsection{Future Directions}

The observed non-inferiority of automatic ROIs motivates
several extensions.  Training the segmentation model directly
on the evaluation cohort could further close the residual
accuracy gap.  End-to-end fine-tuning where the segmentation
network and VLM are jointly optimized---though challenging
due to the VLMs' size---represents a longer-term goal.
Incorporating uncertainty quantification at the segmentation
level may help flag cases where automatic ROIs are
unreliable, allowing selective human review.  Finally,
prospective validation on an independent multicenter cohort
is needed to establish clinical generalizability.

% ============================================================
\section{Conclusion}
\label{sec:conclusion}
% ============================================================

We have shown that UNet+LTAE automatic kidney segmentation
can effectively replace physician-drawn ROIs in a VLM-based
diagnostic pipeline for 99mTc-EC diuretic renography, with
no statistically significant loss of accuracy.  The ablation
experiments confirm that TAC curves are the dominant
diagnostic information source, while image features provide
independent complementary value.  The two-stage VLM pipeline
offers an explainable, structured diagnostic workflow that
achieves clinically promising accuracy (80.77\% with
physician ROIs).  These findings support the feasibility of
fully automated renography interpretation and identify
target areas---particularly mixed obstruction classification
and segmentation model generalization---for further
improvement.

% ============================================================
\section*{Declarations}
% ============================================================

\noindent\textbf{Funding:} This work was supported by [funding
information].
\par\vspace{0.5em}
\noindent\textbf{Conflicts of interest:} The authors declare
no competing interests.
\par\vspace{0.5em}
\noindent\textbf{Ethics approval:} This retrospective study
was approved by the Institutional Review Board of XXX
Hospital (approval number XXX), and the requirement for
informed consent was waived.
\par\vspace{0.5em}
\noindent\textbf{Data availability:} The datasets generated
during this study are available from the corresponding author
on reasonable request.

% ============================================================
% References
% ============================================================
\begin{thebibliography}{99}

\bibitem{oneill}
O'Reilly P, Aurell M, Britton K, et al.
Consensus on diuresis renography for investigating the
dilated upper urinary tract.
J Nucl Med. 1996;37(11):1872--1876.

\bibitem{rossleigh}
Rossleigh MA, Farnsworth RH, Leighton DM, et al.
Technetium-99m dimercaptosuccinic acid scintigraphy studies
of renal function in children with obstructive uropathy.
J Nucl Med. 1993;34(4):614--618.

\bibitem{gordon}
Gordon I, Piepsz A, Sixt R.
Guidelines for standard and diuretic renography in children.
Eur J Nucl Med Mol Imaging. 2011;38(6):1175--1188.

\bibitem{taylor}
Taylor AT, Blaufox MD, De Palma D, et al.
Guidance document for structured reporting of diuresis
renography.
J Nucl Med. 2012;53(5):828--832.

\bibitem{degaard}
DeGrado TR, Bhatt AA, Bhatt S, et al.
Interobserver variability in region of interest selection for
renographic split function determination.
J Nucl Med Technol. 2008;36(1):27--31.

\bibitem{ronneberger}
Ronneberger O, Fischer P, Brox T.
U-Net: convolutional networks for biomedical image
segmentation.
In: MICCAI 2015. LNCS, vol~9351. Springer; 2015:234--241.

\bibitem{isensee}
Isensee F, Jaeger PF, Kohl SAA, et al.
nnU-Net: a self-configuring method for deep learning-based
biomedical image segmentation.
Nat Methods. 2021;18(2):203--211.

\bibitem{openai}
OpenAI.
GPT-4 technical report.
arXiv:2303.08774. 2023.

\bibitem{gpt4}
Achiam J, Adler S, Agarwal S, et al.
GPT-4 technical report.
arXiv:2303.08774. 2024.

\bibitem{wu}
Wu JY, Hsu CC, Chen YC, et al.
ChatGPT-based chest X-ray interpretation: performance and
limitations.
Radiology. 2023;309(1):e231450.

\bibitem{luo}
Luo X, Li Q, Niu Y, et al.
Vision-language models for retinal image analysis: a
systematic review.
Ophthalmol Sci. 2024;4(3):100452.

\bibitem{sibille}
Sibille L, Champsaur D, Duchatelle V, et al.
Natural language processing for nuclear medicine reporting:
current status and future directions.
Eur J Nucl Med Mol Imaging. 2023;50(8):2267--2275.

\bibitem{werner}
Werner RA, Nakajo M, Javadi MS, et al.
Interpretation of nuclear medicine reports using large
language models.
J Nucl Med. 2024;65(2):298--304.

\bibitem{official}
Taylor AT, Blaufox MD, De Palma D, et al.
SNMMI procedure standard/EANM practice guideline for
diuretic renography in adults.
J Nucl Med Technol. 2024;52(1):3--11.

\end{thebibliography}

\end{document}
